\begin{figure}[!ht]
		\centering
		\begin{tikzpicture}[>=Stealth, line width=0.8pt]
			
			% --- coordinate dei nodi ---
			\coordinate (A) at (0,4);
			\coordinate (B) at (3,4);
			\coordinate (C) at (5,4);
			\coordinate (D) at (3,2);
			\coordinate (E) at (5,2);
			\coordinate (F) at (3,0);
			\coordinate (G) at (5,0);
			\coordinate (H) at (8,0);
			
			% --- C_1 = ABCE: traversa superiore A-B-C + asta verticale C-E ---
			\draw[line width=1.5pt] (A) -- (C);
			\draw[line width=1.5pt] (C) -- ($(E)+(0,0.50)$);
			
			% --- C_2 = HGED: H-G + G-E + E-D ---
			\draw[line width=1.5pt] (H) -- (G);
			\draw[line width=1.5pt] (E) -- (G);
			\draw[line width=1.5pt] (E) -- ($(D)+(0.50,0)$);
			
			
			% --- C_3 = BDFG: B-D + D-F + F-G ---
			\draw[line width=1.5pt] (B) -- (D);
			\draw[line width=1.5pt] (D) -- (F);
			\draw[line width=1.5pt] (F) -- ($(G)+(-0.50,0)$);
			
			% =====================================================
			% VINCOLI ESTERNI
			% =====================================================
			
			% --- incastro in A (in alto a sinistra, parete a sinistra) ---
			\draw[thick] (0,3.62) -- (0,4.38);
			\foreach \y in {-0.30,-0.18,-0.06,0.06,0.18,0.30}{\draw[thin] (0,4+\y) -- (-0.12,4+\y-0.10);}
			
			% --- incastro in H (in basso a destra, parete a destra) ---
			\draw[thick] (8,-0.38) -- (8,0.38);
			\foreach \y in {-0.30,-0.18,-0.06,0.06,0.18,0.30}{\draw[thin] (8,\y) -- (8.12,\y-0.10);	}
			
			% =====================================================
			% CONNESSIONI INTERNE
			% =====================================================
			
			% --- cerniera interna in B ---
			\fill[white] ($(B)+(-0.08,0)$) arc(180:360:0.08) -- cycle;
			\draw[thick] ($(B)+(-0.08,0)$) arc(180:360:0.08);
			
			%\draw[fill=white, thick] (B) circle (3pt);
			
			% =====================================================
			% CARRELLO INTERNO in D (asse orizzontale)
			% =====================================================
			
			% rotelle: più piccole (raggio 0.06), tangenti all'asta verticale di C_3
			\draw[thick] ($(D)+(0.09,-0.10)$) circle (0.08);
			\draw[thick] ($(D)+(0.09, 0.10)$) circle (0.08);
			
			% triangolo: base verticale a x = D_x + 0.20 (vertici interni a y = ±0.25),
			%            base prolungata fino a y = ±0.35
			\draw[thick] ($(D)+(0.20,-0.35)$) -- ($(D)+(0.20,0.35)$);
			\draw[thick] ($(D)+(0.20,-0.25)$) -- ($(D)+(0.50,0)$);
			\draw[thick] ($(D)+(0.20, 0.25)$) -- ($(D)+(0.50,0)$);
			
			% pallino bianco al vertice del triangolo (perno della connessione con C_2)
			\node[circle, fill=white, draw=black, inner sep=0pt,	minimum size=4pt, line width=0.7pt] at ($(D)+(0.50,0)$) {};
			
			
			% =====================================================
			% CARRELLO INTERNO in E (asse verticale)
			% =====================================================
			
			% prolungamento del tratto D-E di C_2 oltre E, per accogliere la base del carrello
			\draw[line width=1.5pt] (E) -- ($(E)+(0.20,0)$);
			
			% rotelle: tangenti al tratto orizzontale di C_2 a y = E_y, centri a y = E_y + 0.09
			\draw[thick] ($(E)+(-0.10,0.09)$) circle (0.08);
			\draw[thick] ($(E)+( 0.10,0.09)$) circle (0.08);
			
			% triangolo del carrello: base orizzontale a y = E_y + 0.20, vertice in alto a y = E_y + 0.50
			\draw[thick] ($(E)+(-0.35,0.20)$) -- ($(E)+(0.35,0.20)$);
			\draw[thick] ($(E)+(-0.25,0.20)$) -- ($(E)+(0,0.50)$);
			\draw[thick] ($(E)+( 0.25,0.20)$) -- ($(E)+(0,0.50)$);
			
			% pallino bianco al vertice del triangolo
			\node[circle, fill=white, draw=black, inner sep=0pt,
			minimum size=4pt, line width=0.7pt] at ($(E)+(0,0.50)$) {};
			
			% =====================================================
			% CARRELLO INTERNO in G (asse orizzontale)
			% =====================================================
			
			% prolungamento del tratto E-G di C_2 oltre G, per accogliere la base del carrello
			\draw[line width=1.5pt] (E) -- ($(G)+(0,-0.20)$);
			
			% rotelle: tangenti all'asta verticale di C_2 (E-G), a sinistra
			\draw[thick] ($(G)+(-0.09,-0.10)$) circle (0.08);
			\draw[thick] ($(G)+(-0.09, 0.10)$) circle (0.08);
			
			% triangolo: base verticale a x = G_x - 0.20, vertice a x = G_x - 0.50
			\draw[thick] ($(G)+(-0.20,-0.35)$) -- ($(G)+(-0.20,0.35)$);
			\draw[thick] ($(G)+(-0.20,-0.25)$) -- ($(G)+(-0.50,0)$);
			\draw[thick] ($(G)+(-0.20, 0.25)$) -- ($(G)+(-0.50,0)$);
			
			% pallino bianco al vertice del triangolo (perno della connessione con C_3)
			\node[circle, fill=white, draw=black, inner sep=0pt,
			minimum size=4pt, line width=0.7pt] at ($(G)+(-0.50,0)$) {};
			
			% =====================================================
			% ETICHETTE DEI NODI
			% =====================================================
			\node[font=\scriptsize, above right]       at (A) {$A$};
			\node[font=\scriptsize, above]       at (B) {$B$};
			\node[font=\scriptsize, above]       at (C) {$C$};
			\node[font=\scriptsize, left]        at ($(D)+(-0.12,0)$) {$D$};
			\node[font=\scriptsize, right]       at ($(E)+(0.12,0)$)  {$E$};
			\node[font=\scriptsize, left] at (F) {$F$};
			\node[font=\scriptsize, above right] at (G) {$G$};
			\node[font=\scriptsize, above left]       at (H) {$H$};
			
		\end{tikzpicture}
		\caption{Struttura mista: due mensole ad asse spezzato
			$\mathcal{S}_1 = ABCE$ (incastrata in $A$) e
			$\mathcal{S}_2 = HGED$ (incastrata in $H$),
			collegate da una trave appoggiata $\mathcal{S}_3 = BDFG$
			(cerniera in $B$, carrello in $G$). I carrelli in
			$D$ ed $E$ aggiungono due vincoli semplici allo
			schema isostatico di riferimento: $\nu = -2$,
			iperstatica di grado~2.}
		\label{fig:decomposizione_riduzione}
	\end{figure}
